\section{Why an inner copy of the surface fails}
\label{sec:skeleton}

The layout of the two previous sections places the sources on a smaller copy of the surface. For the cube and the almond this layout does not converge.

\emph{Cube.} For the unit cube at $k = 2\pi$ we placed between $96$ and $600$ electric dipoles on the faces of a concentric inner cube. The error of the scattered field stays between $21\%$ and $30\%$ and does not decrease as sources are added (\autoref{sec:cube}).

\emph{Almond.} For the NASA almond at $1.19$\,GHz, our first layouts on an inner surface were unstable and failed a test with a known exterior solution. A confocal inner surface that resolves the rim gave a stable fit, but its boundary residual remained at $22\%$ with $1243$ sources.

Both failures have the same cause, and it is the rule quoted in the introduction: the sources must enclose the singularities of the continued field. We now locate these singularities for both shapes.

\subsection{Convex polyhedra}

For a polyhedron the continuation can be written down explicitly. Let $\Omega$ be a convex polyhedron. For each face $F$ let $C_F$ be the set of points of $\Omega$ that are closer to the plane of $F$ than to the plane of any other face. The cells $C_F$ fill $\Omega$, and their common walls form the interior medial surface $\mathcal{M}$ of the polyhedron, which we call its \emph{skeleton}. For the cube the skeleton consists of twelve triangles. Each triangle joins an edge to the centre and lies in the plane that bisects the two faces meeting at that edge.

\begin{theorem}
\label{thm:skeleton}
Let $\Omega$ be a convex perfectly conducting polyhedron with skeleton $\mathcal{M}$, and let $\Einc$ be an entire solution of Maxwell's equations. Then the scattered field outside $\Omega$ is the field radiated by an electric surface current $\bJ$ and a magnetic surface current $\bK$ on $\mathcal{M}$.
\end{theorem}

\begin{proof}
Let $(\bE^{t},\bH^{t})$ be the total field outside $\Omega$. Fix a face $F$ with unit normal $\bn$, let $\bx^*$ be the mirror image of $\bx$ in the plane of $F$, and let $P = I - 2\bn\bn^{\top}$. The reflected field
\begin{equation}
  \bE_F(\bx) = -P\,\bE^{t}(\bx^*), \qquad \bH_F(\bx) = P\,\bH^{t}(\bx^*)
\end{equation}
solves Maxwell's equations wherever $\bx^*$ lies outside $\Omega$. Since $\Omega$ is convex, its mirror image lies outside $\Omega$, so $(\bE_F,\bH_F)$ is defined in all of $\Omega$. On $F$ the total field satisfies $\bn\times\bE^{t} = 0$ and therefore $\bn\cdot\bH^{t} = 0$. Hence the reflected field has the same tangential traces on $F$ as the total field, and it continues the total field analytically across the open face. This is the image principle read backwards.

Define $\mathbf{V} = \bE^{t}$ outside $\Omega$ and $\mathbf{V} = \bE_F$ in each cell $C_F$. Then $\mathbf{V}$ solves Maxwell's equations away from $\mathcal{M}$, and $\mathbf{U} = \mathbf{V} - \Einc$ equals the scattered field outside $\Omega$ and satisfies the radiation condition. We apply the Stratton--Chu representation formula \cite{colton2013} to $\mathbf{U}$ on the two sides of $\mathcal{M}$. Small tubes around the edges and small balls around the vertices contribute nothing in the limit, because the field grows at most like $\rho^{-1/3}$ at distance $\rho$ from an edge and like $r^{\nu-1}$ with $\nu > 0$ at distance $r$ from a vertex \cite{meixner1972,costabel2000}, and the reflected fields inherit the same growth. What remains is the field of the currents
\begin{equation}
  \bJ = \bn_{\mathcal{M}}\times[\bH], \qquad \bK = -\bn_{\mathcal{M}}\times[\bE]
  \label{eq:jumps}
\end{equation}
on $\mathcal{M}$, where $[\,\cdot\,]$ denotes the jump between the continuations from the two neighbouring cells.
\end{proof}

The theorem has three consequences for the method.
\begin{enumerate}
  \item Sources on the skeleton are sufficient, and nothing closer to the faces is needed.
  \item The skeleton contains the edges. A shrunk copy of the surface does not reach them and cannot represent the field. This is the stall at $21$--$30\%$.
  \item Both currents in \eqref{eq:jumps} are non-zero in general, so both electric and magnetic dipoles are needed. With electric dipoles alone, our edge sources stalled at an error of $5\%$.
\end{enumerate}
The densities $\bJ$ and $\bK$ are smooth on each wall of the skeleton and become singular only towards the edges and vertices. This is where the sources must cluster.

\begin{remark}
The statement has a relative in potential theory. Gustafsson \cite{gustafsson1998} showed that a convex polyhedron of uniform density has the same exterior Newtonian potential as a measure supported on finitely many planes inside it, which reaches the boundary only at edges and corners. \autoref{thm:skeleton} is the corresponding statement for scattering by a conductor. To our knowledge it has not been stated in this form. The proof uses convexity. For a non-convex polyhedron the mirror image of the body can re-enter the body, and the statement does not hold as it stands.
\end{remark}

\subsection{Smooth thin bodies}

For a smooth surface there is no reflection formula, but the location of the singularities is known in simple cases. For an ellipse in two dimensions the continued field is singular at the two foci \cite{millar1971,barnett2008}. Every cross-section of the almond is an ellipse with axis ratio $3:1$, whose foci lie at the fraction $\sqrt{1-1/9}\approx0.943$ of the semi-major axis. We therefore expect the singularities of the continued field of the almond on the curve traced by these foci in the mid-plane. This \emph{focal curve} runs within $6\%$ of the rim. A shrunk copy of the surface that stays outside the focal curve would have to lie within $6\%$ of the rim itself, where it is too close to the surface to be useful.

This is the standard continuation argument applied section by section, and it is not a proof for the three-dimensional body. Our evidence is numerical. Sources on the mid-plane that cluster towards the focal curve converge (\autoref{sec:almond}), and sources on surfaces that leave the focal curve outside do not.

The skeleton of the cube and the mid-plane of the almond play the same role. In both cases the sources are placed on a surface inside the body that the continued field cannot cross, and they cluster towards the part of that surface which touches or approaches the boundary.
